\chapter{Introdução}

Este documento apresenta a terceira etapa do desenvolvimento do Sistema de Gestão de Clínica, consolidando a transição entre a análise de requisitos e o design detalhado da solução técnica. Após a definição do escopo, dos requisitos funcionais e não funcionais e do mapeamento de casos de uso nas fases anteriores, este relatório avança para a especificação técnica e estrutural necessária para guiar a implementação segura e eficiente do software.

O conteúdo deste documento está organizado em quatro blocos fundamentais:
\begin{enumerate}
    \item \textbf{Mapeamento Objeto-Relacional:} Tradução do modelo de domínio conceitual para a estrutura física de persistência de dados em um banco de dados relacional.
    \item \textbf{Diagrama de Classes em Nível de Projeto:} Detalhamento das classes de software, incluindo seus atributos, métodos, tipos de dados e visibilidade, aproximando a modelagem do código final.
    \item \textbf{Diagramas Dinâmicos:} Representação do comportamento do sistema e da ordem de troca de mensagens entre os componentes durante a execução dos fluxos principais.
    \item \textbf{Arquitetura do Software:} Modelagem estrutural do sistema em alto nível, detalhando como as preocupações das partes interessadas são resolvidas e justificando as escolhas tecnológicas.
    \item  \textbf{Status dos Casos de Uso:} Definição de quais casos de uso já estão implementados.
\end{enumerate}

Com essa estrutura, o documento serve como um guia técnico definitivo para a equipe de desenvolvimento, minimizando ambiguidades e garantindo o alinhamento com as regras de negócio estabelecidas.